\documentclass[aspectratio=169,12pt]{beamer}
\usepackage[utf8]{inputenc}
\usepackage[T1]{fontenc}
\usepackage{lmodern}
\usepackage[spanish]{babel}
% --- Estilo según los Apunts de Comunicacions -----------------------
% Letra sin serifa y muy legible (DejaVu Sans, de la familia de Verdana).
% Con 12pt en un lienzo de 16x9 cm y la letra al 92 %, el cuerpo equivale a
% unos 24 pt en una diapositiva normal, y \large a unos 28 pt.
\usepackage[scaled=0.92]{DejaVuSans}
\renewcommand*\familydefault{\sfdefault}
\usepackage{ccicons}
\usepackage{tikz}

\usetheme{Madrid}

% Colores tranquilos: texto oscuro sobre fondo claro
\definecolor{azul}{HTML}{2E4A5C}
\definecolor{terra}{HTML}{A0583F}
\definecolor{gris}{HTML}{4B5563}
\definecolor{grisclaro}{HTML}{E5E7EB}
\usecolortheme[named=azul]{structure}
\setbeamercolor{normal text}{fg=black!90,bg=white}
\setbeamercolor{alerted text}{fg=terra}
% Título en color oscuro sobre fondo blanco, sin barra de color
\setbeamercolor{frametitle}{fg=azul,bg=white}
\setbeamertemplate{navigation symbols}{}
% Pie discreto: solo el número de transparencia (útil en el turno de
% preguntas). Sin nombre ni título repetidos en cada página.
\setbeamertemplate{footline}{%
  \hfill{\small\color{gris}\insertframenumber\,/\,\inserttotalframenumber}\hspace{0.4cm}\vspace{0.3cm}}
\setbeamerfont{itemize/enumerate body}{size=\large}
\setbeamerfont{itemize/enumerate subbody}{size=\normalsize}
\setbeamerfont{frametitle}{size=\Large,series=\bfseries}
\setlength{\parskip}{0.2cm}

% Cifra grande destacada
\newcommand{\cifra}[2]{%
  \begin{center}
    {\color{terra}\fontsize{36}{40}\selectfont\bfseries #1}\\[0.2cm]
    {\large #2}
  \end{center}}

% Portadilla de cada bloque: {número}{título}{persona}
\newcommand{\bloque}[3]{%
  \begin{frame}[plain]
    \vfill
    \begin{center}
      {\color{terra}\large\bfseries Bloque #1}\\[0.3cm]
      {\color{azul}\LARGE\bfseries #2\par}\vspace{0.5cm}
      {\large #3}
    \end{center}
    \vfill
  \end{frame}}

% Fuente al pie de la transparencia
\newcommand{\fuente}[1]{\vfill{\small\color{gris}#1\par}}

\title[Software: control o participación]{Software como elemento de control o participación de la ciudadanía}
\subtitle{Software libre y desarrollo social}
\author[Grupo SLDS]{Silvia, Ossama, Mar, Héctor y Ona}
\institute{UPC}
\date{Noviembre de 2026\\[0.3cm]{\small\ccbysa\ Contenido bajo licencia CC BY-SA 4.0}}

\begin{document}

% ============================================================
% PORTADA DEL GRUPO
% ============================================================
\begin{frame}
  \titlepage
\end{frame}

% ============================================================
% PRESENTACIÓN
% ============================================================
\section{Presentación}
\begin{frame}{Presentación}
  \vspace{0.2cm}
  {\large\alert{El software decide qué se vigila, qué se publica, cómo se vota y quién accede al Estado.}}

  Solo el software libre permite a la ciudadanía auditar esas decisiones.
  \vspace{0.3cm}
  \begin{enumerate}
    \item Reconocimiento facial y cámaras con IA \hfill \textit{Silvia}
    \item Control cotidiano: del crédito social al altavoz de casa \hfill \textit{Ossama}
    \item Espacio público digital: desinformación y censura \hfill \textit{Mar}
    \item Votación electrónica \hfill \textit{Héctor}
    \item Administración electrónica \hfill \textit{Ona}
  \end{enumerate}
\end{frame}

% ============================================================
% BLOQUE 1 · SILVIA
% ============================================================
\section{Reconocimiento facial}
\bloque{1}{Reconocimiento facial\\y cámaras con IA}{Silvia}

% 2.2 · Cómo identifica una cámara con IA
\begin{frame}{Cómo identifica una cámara con IA}
  \begin{center}
  \begin{tikzpicture}[
      caja/.style={draw=azul,very thick,rounded corners=4pt,minimum width=3.6cm,minimum height=1.5cm,align=center,font=\large},
      flecha/.style={->,very thick,draw=gris}]
    \node[caja] (a) at (0,0) {Captura\\del rostro};
    \node[caja] (b) at (4.7,0) {Vector\\biométrico};
    \node[caja] (c) at (9.4,0) {Comparación\\con una base};
    \draw[flecha] (a) -- (b);
    \draw[flecha] (b) -- (c);
  \end{tikzpicture}
  \end{center}
  \vspace{0.3cm}
  \begin{columns}[T]
    \column{0.45\textwidth}
      \begin{center}
        {\color{azul}\fontsize{30}{34}\selectfont\bfseries 1:1}\\[0.1cm]
        {\large\bfseries Verificación}\\
        {\large «¿Eres quien dices ser?»}
      \end{center}
    \column{0.45\textwidth}
      \begin{center}
        {\color{terra}\fontsize{30}{34}\selectfont\bfseries 1:N}\\[0.1cm]
        {\large\bfseries Identificación masiva}\\
        {\large «¿Quién eres?»}
      \end{center}
  \end{columns}
\end{frame}

% 2.3 · Sesgos
\begin{frame}{El mito de la neutralidad algorítmica}
  {\large Error al clasificar el género (Gender Shades, 2018)}
  \begin{center}
  \begin{tikzpicture}[x=0.2cm,y=1cm,font=\large]
    % 34,7 % ocupa 6,94 cm
    \foreach \y/\emp/\a/\b/\ta/\tb in {0/IBM/34.7/0.3/{34,7}/{0,3}, -1.25/Face++/34.5/0.8/{34,5}/{0,8}, -2.5/Microsoft/20.8/0.05/{20,8}/{0,0}}{
      \node[anchor=east,font=\large\bfseries] at (-1,\y-0.22) {\emp};
      \fill[terra] (0,\y) rectangle (\a,\y-0.38);
      \node[anchor=west,font=\bfseries] at (\a+0.6,\y-0.19) {\ta\,\%};
      \fill[gris] (0,\y-0.45) rectangle (\b,\y-0.83);
      \node[anchor=west,font=\bfseries] at (\b+0.6,\y-0.64) {\tb\,\%};
    }
  \end{tikzpicture}\\[0.2cm]
  {\large\textcolor{terra}{\rule{0.35cm}{0.35cm}}~Mujeres de piel oscura\qquad\textcolor{gris}{\rule{0.35cm}{0.35cm}}~Hombres de piel clara}
  \end{center}
  {\large NIST: falsos positivos de \alert{\textbf{10 a 100 veces}} más según demografía}
  \fuente{Buolamwini y Gebru, PMLR 81 (2018) · NIST IR 8280 (2019)}
\end{frame}

% 2.4 · Marco legal
\begin{frame}{El freno europeo: EU AI Act}
  \begin{itemize}
    \item Biometría remota en tiempo real: \alert{\textbf{prohibida}}*
    \item Bases faciales por \textit{scraping} masivo: prohibidas
    \item Identificación diferida: sistema de alto riesgo
    \item ¿Y una matrícula? \textbf{No es biometría}
  \end{itemize}
  \fuente{* Uso policial, salvo 3 excepciones · Reglamento (UE) 2024/1689}
\end{frame}

% 2.5 · Flock: tecnología
% (Si conseguís una foto libre de una cámara Flock, de 1024x768 o más y
% con licencia CC BY / CC BY-SA, ponedla en una columna a la derecha.)
\begin{frame}{Flock Safety: vigilancia sin rostro}
  \begin{itemize}
    \item \alert{\textbf{+120.000}} cámaras en 49 estados
    \item Sin reconocimiento facial: «huella del vehículo»
    \item Matrícula, color, marca, baca, pegatinas
  \end{itemize}
  \fuente{Sen. Hawley (26/08/2026) · Robinson v. Commonwealth (Virginia, 2026)}
\end{frame}

% 2.6 · Flock: escala
% CORREGIDO 23/11/2026: en Robinson se buscó por marca y modelo (las llantas
% sirvieron después para reconocer el coche). Se quita «Red compartida entre
% agencias de todo el país» porque ninguna fuente del pie lo respalda.
\begin{frame}{De leer matrículas a seguir personas}
  \begin{itemize}
    \item Busca coches \alert{\textbf{sin matrícula}}: marca y modelo
    \item Nova: de la matrícula a la persona
    \item OS Investigate: la IA deduce tus «asociados»
  \end{itemize}
  \fuente{Robinson v. Commonwealth · 404 Media (2025) · WIRED (2026)}
\end{frame}

% 2.7 · Controversias
\begin{frame}{Abusos sin orden judicial}
  \begin{itemize}
    \item Sin orden judicial: los tribunales lo avalan, por ahora
    \item Al menos 50 agentes acusados de uso indebido
    \item 26 de ellos, por espiar a parejas o exparejas
    \item Un agente: \alert{\textbf{2.048}} búsquedas sobre su víctima
  \end{itemize}
  \fuente{Robinson v. Commonwealth · Washington Post (04/08/2026) · CNN}
\end{frame}

% 2.8 · Caso inmigración
\begin{frame}{La puerta trasera de la red}
  \begin{columns}[T]
    \column{0.31\textwidth}
      \begin{center}
        {\color{azul}\fontsize{44}{48}\selectfont\bfseries 8}\\[0.2cm]
        {\large\bfseries Puerta delantera}\\[0.1cm]
        {\large Acceso directo dado a la Patrulla Fronteriza}
      \end{center}
    \column{0.31\textwidth}
      \begin{center}
        {\color{terra}\fontsize{44}{48}\selectfont\bfseries 10}\\[0.2cm]
        {\large\bfseries Puerta trasera}\\[0.1cm]
        {\large Redes consultadas por la Patrulla sin permiso}
      \end{center}
    \column{0.31\textwidth}
      \begin{center}
        {\color{azul}\fontsize{44}{48}\selectfont\bfseries 89}\\[0.2cm]
        {\large\bfseries Puerta lateral}\\[0.1cm]
        {\large Redes en una búsqueda local con motivo «ICE»}
      \end{center}
  \end{columns}
  \fuente{UW Center for Human Rights, «Leaving the Door Wide Open», 2025}
\end{frame}

% 2.9 · Reacción social
% CORREGIDO 23/11/2026: 303 no se pudo confirmar; última cifra verificable
% del Institute for Justice, 247 (The Marshall Project, 03/10/2026).
\begin{frame}{La presión ciudadana frena la red}
  \begin{itemize}
    \item \alert{\textbf{247}} gobiernos locales abandonan sus lectores de matrículas
    \item Denver y Mountain View rompen con Flock
    \item Amazon Ring cancela su alianza (feb. 2026)
    \item El Senado de EE.\,UU. abre una investigación
  \end{itemize}
  \fuente{Institute for Justice, vía The Marshall Project (03/10/2026) · Ring (12/02/2026) · Sen. Hawley}
\end{frame}

% 2.10 · Alternativa libre
% (Si añadís una captura del mapa de maps.deflock.org u OpenStreetMap,
% indicad «Datos © OpenStreetMap contributors, ODbL».)
% CORREGIDO 23/11/2026: Pew cuenta las cámaras en OpenStreetMap, no «de DeFlock».
\begin{frame}{DeFlock: auditar con software libre}
  \begin{itemize}
    \item DeFlock (MIT): mapa ciudadano en \textbf{OpenStreetMap}
    \item \alert{\textbf{109.577}} cámaras Flock en OpenStreetMap
    \item Frigate (MIT): matrículas en tu propio servidor
    \item Local $\neq$ inocuo: cambia quién controla
  \end{itemize}
  \fuente{deflock.org · Pew Research Center (09/11/2026) · Frigate NVR}
\end{frame}

% 2.11 · Conclusión del bloque
\begin{frame}{En la calle, el código ya es ley}
  \vspace{0.3cm}
  {\Large Infraestructura \alert{\textbf{privada}} y opaca: vigilancia sin control democrático.\par}
  \vspace{0.8cm}
  {\Large Código y datos abiertos: de promesas a garantías verificables.\par}
\end{frame}

% ============================================================
% BLOQUE 2 · OSSAMA
% ============================================================
\section{Control cotidiano}
\bloque{2}{Control cotidiano: del crédito\\social al altavoz de casa}{Ossama}

% ============================================================
% BLOQUE 3 · MAR
% ============================================================
\bloque{3}{Espacio público digital:\\desinformación y censura}{Mar}

\section{Introducción}
\begin{frame}{Tres decisiones que hoy toma el software}
  \begin{itemize}
    \item \textbf{Qué vemos:} algoritmos que ordenan el debate público.
    \item \textbf{Qué creemos:} microsegmentación y deepfakes.
    \item \textbf{Si podemos conectarnos:} apagones de internet.
  \end{itemize}
  \vspace{0.5cm}
  {\large\alert{Y una respuesta: el software libre.}}
\end{frame}

\begin{frame}{Qué tiene que ver el SW libre?}
  \begin{itemize}
    \item \textbf{UN CÓDIGO QUE CUALQUIERA PUEDE REVISAR Y QUE NADIE PUEDE CONTROLAR NI APAGAR EN SOLITARIO}
  \end{itemize}
  \vspace{0.5cm}
\end{frame}

\section{Desinformación}
\begin{frame}{Cambridge Analytica}
  \cifra{87 millones}{de perfiles recogidos con solo 305.000 tests}
  \begin{itemize}
    \item Facebook dejaba a las apps leer los datos de los amigos.
    \item Se usaron para mensajes políticos a medida (EE.\,UU., 2016).
    \item Multa récord de 5.000 M\$ de la FTC.
  \end{itemize}
\end{frame}

\begin{frame}{Deepfake: el audio falso de Eslovaquia (2023)}
  \begin{itemize}
    \item Dos días antes de votar aparece un audio falso de un candidato.
    \item Cae en el silencio electoral: casi no se puede desmentir.
    \item El candidato, primero en las encuestas, pierde.
  \end{itemize}
  \vspace{0.5cm}
  {\large\alert{Matiz:} no fue la única causa, pero un audio clonado es muy difícil de detectar.}
\end{frame}

\begin{frame}{Algoritmos: una caja negra}
  \begin{itemize}
    \item Optimizan la atención, no la información.
    \item Su código no es público.
    \item Crean cámaras de eco.
    \item La UE (DSA, art.\,38) obliga a ofrecer un feed sin perfilado.
  \end{itemize}
  \vspace{0.5cm}
  {\large\alert{Solo el código abierto permite auditar.}}
\end{frame}

\section{Apagones}
\begin{frame}{Cómo se apaga internet}
  \begin{itemize}
    \item \textbf{BGP:} el país desaparece del mapa de internet (Egipto, 2011).
    \item \textbf{DNS:} se bloquean webs concretas (Turquía, 2014).
    \item \textbf{Ralentizar:} la red funciona, pero no sirve (Irán, 2026).
    \item \textbf{Destruir:} se bombardean cables y antenas (Gaza).
  \end{itemize}
  \cifra{313 apagones}{en 52 países solo en 2025}
\end{frame}

\begin{frame}{Tres casos}
  \begin{columns}[T]
    \column{0.32\textwidth}
      \cifra{85}{apagones en Myanmar (2024)}
    \column{0.32\textwidth}
      \cifra{$\sim$70\,\%}{de 2026 sin internet en Irán}
    \column{0.32\textwidth}
      \cifra{7 días}{de apagón total en Gaza (2024)}
  \end{columns}
  \vspace{0.4cm}
  \begin{center}
    Los apagones acompañan a la represión y a la guerra.
  \end{center}
\end{frame}

\begin{frame}{Censura y desinformación se refuerzan}
  \vspace{0.5cm}
  \begin{center}
    {\Large Con Irán desconectado, el régimen pudo hacer pasar vídeos reales por deepfakes.}\\[1cm]
    \alert{\Large Sin red, nadie puede verificar.}
  \end{center}
\end{frame}

\section{Software libre}
\begin{frame}{Tor y Snowflake: todos podemos ser puente}
  \begin{enumerate}
    \item Una persona censurada no puede conectarse a Tor.
    \item Se conecta a través del navegador de un voluntario.
    \item Desde ahí llega a internet sin censura.
  \end{enumerate}
  \cifra{9,6 millones}{de iraníes al día usaron Psiphon en 2026}
\end{frame}

\begin{frame}{Redes sin centro}
  \begin{columns}[T]
    \column{0.48\textwidth}
      \textbf{Redes mesh} (Briar)
      \begin{itemize}
        \item De móvil a móvil, sin operador.
        \item Funcionan aunque caiga la red.
      \end{itemize}
    \column{0.48\textwidth}
      \textbf{Fediverso} (Mastodon)
      \begin{itemize}
        \item Miles de servidores independientes.
        \item Nadie puede apagarlo entero.
      \end{itemize}
  \end{columns}
  \vspace{0.6cm}
  {\large\alert{Si no hay centro, no hay interruptor.}}
\end{frame}

\section{Conclusión}
\begin{frame}{Conclusión}
  \begin{center}
    {\Large Quien controla el código controla la información.}
  \end{center}
  \vspace{0.3cm}
  \begin{itemize}
    \item Algoritmos abiertos y auditables.
    \item Una red libre que permita verificar.
    \item Herramientas sin interruptor central.
  \end{itemize}
\end{frame}

% ============================================================
% BLOQUE 4 · HÉCTOR
% ============================================================
\section{Votación electrónica}
\bloque{4}{Votación electrónica}{Héctor}

% 4.1 · Apertura y mapa del bloque
\begin{frame}{¿Por qué votar no es como hacer un Bizum?}
  \begin{enumerate}
    \item \textbf{Papel frente a urna digital:} qué cambia y qué está en juego.
    \item \textbf{Experiencia internacional:} qué ha pasado al probarlo.
    \item \textbf{Código cerrado frente a libre:} Voatz, Helios y Belenios.
    \item \textbf{Blockchain:} cómo funciona y dónde están sus límites.
  \end{enumerate}
  \vspace{0.4cm}
  {\large\alert{La ciudadanía debe poder verificar cada paso del proceso electoral.}}
\end{frame}

% 5.2 · Tabla comparativa: papel vs. voto electrónico
\begin{frame}{Voto en papel frente a voto electrónico}
  \begin{center}
  \small
  \renewcommand{\arraystretch}{1.25}
  \begin{tabular}{p{2.4cm} p{4.6cm} p{4.6cm}}
    & {\color{azul}\normalsize\bfseries Voto en papel} & {\color{terra}\normalsize\bfseries Voto electrónico} \\
    \hline
    \textbf{Recuento} & \raggedright Manual, lento y con posible error humano & \raggedright Automático, rápido y menos errores de conteo \tabularnewline
    \hline
    \textbf{Logística} & \raggedright Urnas, papeletas y transporte masivo & \raggedright Mayor accesibilidad (incluso desde el móvil) \tabularnewline
    \hline
    \textbf{Riesgos} & \raggedright Manipulación localizada & \raggedright Ciberataques y fallos a gran escala \tabularnewline
    \hline
    \textbf{Supervisión} & \raggedright Proceso físico visible y comprensible & \raggedright Opaco si no es auditable \tabularnewline
  \end{tabular}
  \end{center}
  \vspace{0.2cm}
  {\large\alert{Objetivo del voto electrónico:} ser \textbf{verificable}, \textbf{secreto} y \textbf{auditable}.}
\end{frame}

% 5.3 · Países que apuestan por el voto tradicional
\begin{frame}{Voto tradicional: España y Alemania}
  \begin{columns}[T]
    \column{0.47\textwidth}
      {\color{azul}\large\bfseries España}\\[0.15cm]
      \begin{itemize}
        \item Voto presencial en papel como norma estatal.
        \item Pruebas piloto en 2004 y 2005.
        \item Sin voto electrónico remoto con validez general.
      \end{itemize}
    \column{0.47\textwidth}
      {\color{terra}\large\bfseries Alemania (2009)}\\[0.15cm]
      \begin{itemize}
        \item El Constitucional anuló las máquinas de 2005.
        \item Ni votantes ni mesas podían verificar el registro.
        \item Exigió supervisión sin conocimientos técnicos.
      \end{itemize}
  \end{columns}
  \vspace{0.3cm}
  {\large\alert{No se prohibió votar electrónicamente, sino no poder comprobar el escrutinio.}}
\end{frame}

% 5.4 · Países que lo han implementado o intentado
\begin{frame}{Países que han probado el voto electrónico}
  \begin{columns}[T]
    \column{0.47\textwidth}
      {\color{azul}\large\bfseries Brasil}\\[-0.15cm]
      \begin{itemize}
        \item Urnas electrónicas desde 2000, siempre \textbf{presenciales}.
      \end{itemize}
      \vspace{0.1cm}
      {\color{azul}\large\bfseries India}\\[-0.15cm]
      \begin{itemize}
        \item Desde 1989; $\sim$1\,M de máquinas en 2004 ($>$670\,M de votantes).
        \item Fallos revelados en 2006.
      \end{itemize}
    \column{0.47\textwidth}
      {\color{terra}\large\bfseries Estonia}\\[-0.15cm]
      \begin{itemize}
        \item Voto online desde 2005 (3,13\,\% en 2007), con \textbf{papel de respaldo}.
        \item Auditorías señalan riesgos.
      \end{itemize}
      \vspace{0.1cm}
      {\color{terra}\large\bfseries Países Bajos}\\[-0.15cm]
      \begin{itemize}
        \item En 2007 volvió al papel tras fallos graves en sus máquinas.
      \end{itemize}
  \end{columns}
  \vspace{0.15cm}
  {\large\alert{Implementar el voto electrónico no es una progresión lineal.}}
\end{frame}

% 5.5 · Ejemplo privativo: Voatz
\begin{frame}{Código cerrado: el caso Voatz}
  \begin{itemize}
    \item App móvil de votación con \textit{blockchain} y \textbf{código privativo}.
    \item Auditoría del \textbf{MIT}: fallos para alterar, impedir o revelar el voto.
    \begin{itemize}
      \item \textit{Malware} o \textit{root} tomaba el control; la red filtraba datos del voto.
      \item Un tercero podía exponer fotos y datos de identidad.
      \item \textit{Blockchain} no protegía frente a ataques al servidor.
    \end{itemize}
    \item \textbf{Trail of Bits} confirmó fallos en auditoría y control de administradores.
  \end{itemize}
  \vspace{0.2cm}
  {\large\alert{Usar \textit{blockchain} no garantiza seguridad: el eslabón más débil es el dispositivo.}}
  \fuente{MIT (2020) · Trail of Bits (2020)}
\end{frame}

% 5.6 · Tabla comparativa: Voatz vs. Helios vs. Belenios
\begin{frame}{Código cerrado frente a software libre}
  \begin{center}
  \small
  \renewcommand{\arraystretch}{1.15}
  \begin{tabular}{p{1.9cm} p{3.3cm} p{3.3cm} p{3.4cm}}
    & {\color{terra}\normalsize\bfseries Voatz} & {\color{azul}\normalsize\bfseries Helios} & {\color{azul}\normalsize\bfseries Belenios} \\
    \hline
    \textbf{Licencia} & \raggedright Privativo (app móvil + \textit{blockchain}) & \raggedright Software libre (web, Ben Adida) & \raggedright Software libre (evolución de Helios) \tabularnewline
    \hline
    \raggedright\textbf{Punto fuerte} & \raggedright Comodidad móvil & \raggedright Cifrado y pruebas verificables & \raggedright Más privacidad y defensa ante \textit{ballot stuffing} \tabularnewline
    \hline
    \textbf{Límite} & \raggedright Opaco; fallos críticos en auditorías & \raggedright Confiar en que la urna no añada votos & \raggedright Depende del dispositivo y credenciales \tabularnewline
  \end{tabular}
  \end{center}
  \vspace{0.1cm}
  {\large\alert{El software libre no elimina fallos, pero permite auditar y reproducir el recuento.}}
\end{frame}

% 5.7 · Blockchain: flujo de 5 pasos + 5 capas de seguridad
\begin{frame}{Blockchain: flujo y capas de seguridad}
  \begin{center}
  \begin{tikzpicture}[
      paso/.style={draw=azul,very thick,rounded corners=4pt,minimum width=2.15cm,minimum height=1.1cm,inner sep=4pt,align=center,font=\scriptsize\bfseries},
      flecha/.style={->,very thick,draw=gris}]
    \node[paso] (p1) at (0,0) {1. Identidad\\digital};
    \node[paso] (p2) at (2.75,0) {2. Cifrado\\en el móvil};
    \node[paso] (p3) at (5.50,0) {3. Registro\\en cadena};
    \node[paso] (p4) at (8.25,0) {4. Descifrado\\y recuento};
    \node[paso] (p5) at (11.00,0) {5. Pruebas\\verificables};
    \draw[flecha] (p1) -- (p2);
    \draw[flecha] (p2) -- (p3);
    \draw[flecha] (p3) -- (p4);
    \draw[flecha] (p4) -- (p5);
  \end{tikzpicture}
  \end{center}
  \vspace{0.2cm}
  {\large\textbf{Cinco piezas de la seguridad real}} {\normalsize (Helios ya verifica sin \textit{blockchain}):}
  \begin{columns}[T]
    \column{0.48\textwidth}
      \begin{itemize}
        \item Criptografía robusta
        \item Claves repartidas entre autoridades
        \item Verificabilidad pública
      \end{itemize}
    \column{0.48\textwidth}
      \begin{itemize}
        \item Auditorías independientes
        \item \alert{\textbf{Papel de respaldo}}
      \end{itemize}
  \end{columns}
\end{frame}

% 5.8 · Límites y conclusión del bloque
\begin{frame}{Dos muros que la tecnología no derriba}
  \begin{columns}[T]
    \column{0.47\textwidth}
      {\color{terra}\large\bfseries Coacción y compra de votos}\\[0.15cm]
      \begin{itemize}
        \item Sin cabina electoral en casa.
        \item Alguien puede ver tu pantalla o exigir tu clave.
        \item Debilidad estructural del voto online.
      \end{itemize}
    \column{0.47\textwidth}
      {\color{terra}\large\bfseries Malware en el dispositivo}\\[0.15cm]
      \begin{itemize}
        \item Cambia el voto \textbf{antes} de cifrarlo y te engaña.
        \item La criptografía no arregla un móvil comprometido.
        \item Riesgos de duplicados y ataques DoS.
      \end{itemize}
  \end{columns}
  \vspace{0.25cm}
  \begin{center}
    {\normalsize\alert{\textbf{Sin software libre y sin supervisión ciudadana, la tecnología no refuerza la democracia: solo la vuelve opaca.}}}
  \end{center}
\end{frame}

% ============================================================
% BLOQUE 5 · ONA
% ============================================================
\section{Administración electrónica}
\bloque{5}{Administración electrónica}{Ona}


\end{document}

